\section{Discussion}

\paragraph{Key Insights.}
Our results demonstrate that versioned document QA requires fundamentally different retrieval mechanisms than standard RAG. The 32-point accuracy gap between VersionRAG and Naive RAG cannot be bridged by simply retrieving more context or using larger models—it stems from the inability to distinguish between version-valid and version-invalid information. Even GraphRAG, despite its sophisticated entity-relationship modeling, achieves only marginal improvements (64\% vs 58\%) because version transitions are not explicit relationships that can be extracted from text. They must be modeled as first-class citizens in the retrieval system.

\paragraph{Practical Deployment Strategy.}
The correlation between model capacity and performance (55\% to 90\% accuracy) suggests a tiered deployment. For simple version lookups and content retrieval where VersionRAG maintains reasonable accuracy (61.6\%) even with Llama 3 8B, smaller models suffice. Complex change detection and multi-version reasoning tasks require larger models for accurate query classification and change interpretation. Organizations can optimize costs by routing queries based on complexity—a strategy our query classification component naturally supports.

\paragraph{Generalizability Beyond Technical Documentation.}
While evaluated on technical documentation, VersionRAG's principles apply to any domain with discrete document versions. Legal documents undergo formal revisions with tracked changes. Scientific papers have preprint versions and post-review updates. Medical guidelines and regulatory compliance documents follow strict versioning for liability reasons. The key requirement is identifiable version markers and structural consistency across versions. Continuous update streams (news articles) or unstructured versioning (social media) would require different approaches, particularly for change detection.

\paragraph{Efficiency-Flexibility Trade-off.}
VersionRAG's 97\% reduction in indexing tokens compared to GraphRAG represents a design philosophy: leveraging domain knowledge about versioning patterns rather than discovering all possible relationships. This specialization enables practical deployment for large document collections while maintaining superior accuracy for version-specific tasks. GraphRAG's flexibility to discover unexpected relationships is valuable for exploratory analysis. Future work could explore hybrid approaches that combine VersionRAG's efficient version tracking with GraphRAG's open-ended relationship discovery for comprehensive document understanding.

\paragraph{Implications for RAG System Design.}
Our work highlights that domain-specific structures in documents—whether versions, sections, or temporal markers—should be explicitly modeled rather than hoping generic retrieval will capture them. The success of purpose-built components (version-aware graph, query classification, change detection) over general-purpose solutions suggests that RAG systems benefit from domain-aware architectures. This challenges the trend toward ever-larger, general-purpose models and retrieval systems, advocating instead for structured approaches that encode domain knowledge into the system architecture.


\section{Conclusion}
We presented VersionRAG, the first retrieval-augmented generation framework specifically designed for versioned documents. By introducing a version-aware graph structure that explicitly models document evolution and changes, VersionRAG achieves 90\% accuracy on version-sensitive questions—a 32-point improvement over standard RAG and 26-point improvement over GraphRAG. Beyond effectiveness, VersionRAG demonstrates remarkable efficiency, requiring 97\% fewer tokens during indexing than GraphRAG while maintaining superior performance.

Our work makes three key contributions to the field. First, we formalize the versioned document QA task and identify three fundamental query types that require distinct retrieval strategies, providing a framework for future research. Second, we develop VersionRAG's novel graph-based architecture that bridges the gap between efficient vector retrieval and precise version-aware filtering. Third, we create VersionQA, the first benchmark for evaluating version-aware document QA systems, comprising 100 manually curated questions across 34 versioned documents.

The implications extend beyond technical documentation. As organizations increasingly rely on LLMs to navigate complex document repositories, the ability to reason about document evolution becomes critical for applications in compliance, debugging, and knowledge management. VersionRAG's efficient indexing and incremental update capabilities make it practical for deployment on continuously evolving document collections.

Future work should address three key areas: (1) expanding evaluation to other domains such as legal and medical documents, (2) developing more sophisticated change detection mechanisms that can capture subtle semantic shifts, and (3) exploring hybrid approaches that combine VersionRAG's version awareness with GraphRAG's open-ended relationship discovery. As LLMs become more capable, we also envision opportunities for learning version patterns from data rather than imposing predefined structures.


\clearpage
